\section{Introduction}
\label{sec:introduction}

% ROLE: opening -- define the task and why both prediction and explanation matter.
Many scientific measurements arrive as families of response curves rather than
independent rows.  A new material or compound is observed at only a few support
conditions, yet the scientist wants both the remaining response curve and a
compact description of how that entity differs from previously measured ones.
We study this support-conditioned setting: given sparse measurements of an
unseen entity, infer entity-specific coordinates from the support alone and use
them to explain its hidden query response.  Thermoelectric figure of merit,
vapor pressure, and crystal heat capacity illustrate the central tension.  A
flexible predictor can interpolate these curves accurately, but a scientific
coordinate must additionally be stable enough to name, compare, and place in a
testable equation.

% ROLE: challenge -- expose latent chart ambiguity as the technical obstacle.
Few-shot function models and implicit neural representations already provide
powerful mechanisms for predicting a new function from a context set
\citep{garnelo2018conditional,xu2020metafun,gondal2021function,dupont2022functa}.
Their internal task vectors, however, are generally coordinates of a learned
chart rather than intrinsic scientific variables.  An invertible rotation,
scaling, or affine change of a latent code can be absorbed by the decoder
without changing a single physical prediction, while changing every individual
coordinate supplied to symbolic regression.  Low dimensionality alone does not
ensure identifiability, and common regularized disentanglement objectives do not
remove the need for explicit assumptions or inductive biases
\citep{locatello2019challenging,khemakhem2020variational}.  Moreover,
canonicalization itself has nontrivial scope and continuity limits
\citep{ma2024canonicalization,dym2024equivariant}.  Consequently, a predictive
raw latent vector cannot be named as a physical quantity merely because it
supports accurate reconstruction.

% ROLE: gap -- distinguish our interface from existing equation discovery.
Scientific equation discovery faces the complementary problem.  Sparse
dynamics, neural-to-symbolic extraction, and parametric equation models can
recover compact relations from suitable coordinates or learned components
\citep{champion2019data,cranmer2020discovering,zhang2022parametric,podina2023universal}.
Yet when an equation coefficient is inferred for an unseen entity, the input
coordinate must be determined from that entity's support and must not depend on
an arbitrary neural chart.  Existing ingredients therefore leave an interface
question unanswered: what representation-level object should be interpreted
when different latent parameterizations predict the same scientific response?

% ROLE: method -- state the invariant object and the concrete pipeline.
We answer this question with \emph{auditable basis-relative response
coordinates}.  We fix a named structure a priori or select it using
development entities only; when refinements are allowed, a bounded theory-guided grammar, grouped
support/query scoring, and a fixed tie break select among them.  Its basis and
probe grid are then sealed for external evaluation.  We evaluate the decoder on
those probes and project its response
onto that basis.  Two representations that agree on the declared probe response
then have exactly the same coefficients within that frozen specification, regardless of their latent
dimension, redundancy, or chart.  For a new entity, stable response-metric
  Gauss--Newton calibration uses only its support; direct structure re-estimation
provides the auditable endpoint.  We additionally formulate a decoder-prior
diagnostic, GIRD, whose singular-direction analysis predicts benefit primarily
when support is underidentified.  Our controlled experiment selects its prior
weight inside each predeclared support regime; it does not instantiate a
per-entity rank router.  This separates three
questions that are often conflated: whether support contains entity state,
whether that state admits named response coordinates, and whether a learned
decoder prior adds information beyond the support.

% ROLE: evidence -- lead with external expression transfer and retain the strongest baseline.
The resulting interpretation target is deliberately operational.  An expression
succeeds when an unseen entity's support-only coefficients yield finite
physical-unit query predictions with pooled $R^2\geq0.85$ and its compact named
terms offer a scientifically suggestive stage-wise reading.  It need not recover
the initial raw latent code, a preselected physical variable, or a unique
microscopic law.  We report coefficient stability and falsifiable follow-up tests
as evidence for that reading rather than treating a particular latent chart as
the target.  Under a frozen temporal protocol, the
quadratic response chart
$\mathrm{ZT}(T)=c_0+c_1\tau+c_2\tau^2$ reaches $R^2=0.9888$ on 30 future
materials, and a reference-pressure/enthalpy/curvature expression reaches
$R^2=0.9996$ on 84 future vapor-pressure compounds.  The coordinates are stable
under support offsets and suggest direct tests of temperature turnover and
effective vaporization-enthalpy trends.  Strong local interpolation remains
competitive or better---log-pressure PCHIP is more accurate for vapor pressure, and the ZT
advantage over kNN is not significant---so our claim is an auditable scientific
coordinate interface rather than universal predictive superiority.

% ROLE: contributions -- every item maps to evidence already present or to the frozen branch rule.
Our contributions are:
\begin{itemize}
  \item We formulate interpretation on equivalence classes of fixed physical
  probe responses and show that named basis coefficients are invariant to any
  response-equivalent representation; we separately characterize the narrower
  affine/full-rank conditions under which support calibration is equivariant.
  \item We provide a support-only, basis-relative inference protocol with explicit
  development-only structure selection, external freezing, rank, conditioning,
  gauge, leakage, coefficient-stability, and entity-tail audits, together with a
  rank-aware test of whether a learned decoder prior adds value beyond direct
  structure re-estimation.
  \item We evaluate compact equations under entity- and time-disjoint real-data
  protocols, reporting uncertainty and strong support-aware baselines separately
  from expression fidelity.  Two domains already provide one-shot temporal
  confirmation; the Crystal heat-capacity stress test follows a frozen
  result-to-narrative rule.
\end{itemize}

% BRANCH RULE: if the terminal learned-prior development and temporal gates fail,
% keep GIRD as a diagnostic in the second contribution and do not call it a
% headline algorithm in the abstract or title.  No partial result may change this text.
